% ============================
% CV-2026 Updated Memo (Alexei Kornaev, Kirill Yakovlev)
% ============================
\documentclass[a4paper,11pt]{texMemo}
\usepackage{graphicx, lipsum}
\usepackage{hyperref}
\usepackage{xcolor}

\usepackage[english]{babel}
\usepackage[square,numbers]{natbib}
\bibliographystyle{abbrvnat}

\memoto{M25-RO-01 (master students)}
\memofrom{Alexei Kornaev, Kirill Yakovlev}
\memosubject{Course on Computer Vision (CV-2026, Semester 2)}
\memodate{\today}
\logo{\includegraphics[width=0.3\textwidth]{logo.png}}

\begin{document}
\maketitle

% ============================
\section{Lectures}
% ============================
\begin{enumerate}

% ------------------------------------------------------------
\item \textbf{Intro to Computer Vision}
\citep{HuggingFaceCVcourse,szeliski2022computer, merkulov2022optimSeminars, prince2023understanding,CVcourseCSCI1430,gitCV-2026, LightningColab}, 
emphasis on \href{https://huggingface.co/learn/computer-vision-course/en/unit1/chapter1/motivation}{Unit 1} \citep{HuggingFaceCVcourse} 
\begin{itemize}
    \item \textbf{Theory:}
    What CV solves today: recognition, localization, segmentation, tracking, 3D understanding, and multimodal reasoning.
    The modern stack: \emph{data \textrightarrow\ augmentation \textrightarrow\ backbone \textrightarrow\ head \textrightarrow\ evaluation \textrightarrow\ deployment}.
    Image formation + photometric distortions (exposure, blur, noise), and why ``dataset shift'' is the default.
    Deterministic CV (filtering/edges/geometry) as baselines and as debugging tools for DL.
    Metrics and pitfalls: accuracy vs F1 vs mAP, calibration, robustness, and reproducibility.
    Preview of the semester arc: \emph{CNNs/ViTs} $\rightarrow$ \emph{promptable/open-vocabulary CV} $\rightarrow$ \emph{VLMs} $\rightarrow$ \emph{VLAs/VLAMs}.
    \item \textbf{Skills :}
    OpenCV as a fast baseline and debugging tool (preprocessing, geometry, tracking primitives).
    PyTorch + Lightning (clean training loops) + ClearML (experiment tracking).
    A first ``CV systems'' lab mindset: set a baseline, log metrics, visualize failure modes, write a short experiment report.
\end{itemize}

% ------------------------------------------------------------
\item \textbf{Convolutional Neural Networks}
\citep{prince2023understanding, HuggingFaceCVcourse, YandexHandbookML,Goodfellow-et-al-2016,howard2020deep, kiranyaz20211d, li2021survey}, emphasis on \href{https://udlbook.github.io/udlbook/}{UDL} book, \href{https://education.yandex.ru/handbook/ml/article/svyortochnye-nejroseti}{Yandex handbook} \cite{prince2023understanding,YandexHandbookML} and \href{https://github.com/fastai/fastbook/blob/master/13_convolutions.ipynb}{hands-on coding chapter 13} \cite{howard2020deep} 
\begin{itemize}
    \item \textbf{Theory :}
    Convolutions as \emph{inductive bias} (locality + translation equivariance), and where CNNs still win (efficiency, small data, edge devices).
    Modern training recipes: normalization (BN/LN), augmentation, regularization, LR schedules, mixed precision.
    Architecture patterns: residual connections, bottlenecks, depthwise separable convs, inverted residuals.
    Minimal but relevant zoo: AlexNet/VGG (history), ResNet (baseline), EfficientNet/MobileNet (efficiency), ConvNeXt (modernization of CNNs).
    When to choose CNN vs ViT, and hybrid backbones.
    \item \textbf{Skills :}
    Implement + debug CNN training (CIFAR-10 or a small robotics-themed dataset).
    Transfer learning and fine-tuning (ResNet/ConvNeXt) with proper layer freezing/unfreezing.
    Build a small ``edge-ready'' model (MobileNet-style) and measure speed/accuracy trade-offs.
\end{itemize}

% ------------------------------------------------------------
\item \textbf{Visual transformers}
\citep{prince2023understanding, vaswani2023attentionneed, dosovitskiy2021imageworth16x16words}, emphasis on \href{https://arxiv.org/abs/2010.11929}{paper on ViTs} \citep{dosovitskiy2021imageworth16x16words}, and \href{https://udlbook.github.io/udlbook/}{UDL} book \citep{prince2023understanding}
\begin{itemize}
    \item \textbf{Theory:}
    Transformer essentials for vision: patching, tokenization, positional encoding, self-attention, MLP blocks.
    Why ViTs scale well and why they may be fragile without good recipes/data.
    ``Transformer family'' overview (no overload): ViT (baseline), hierarchical attention idea (e.g., Swin-style), and why many detectors/segmenters are now transformer-based.
    Attention maps as an interpretability/debugging instrument (with caveats).
    \item \textbf{Skills:}
    Fine-tune a ViT from Hugging Face on a moderate dataset; compare to a CNN baseline.
    Measure compute/memory, and do a short ablation: resolution, augmentation, freezing strategy.
\end{itemize}

% ------------------------------------------------------------
\item \textbf{Generative Models}
\citep{YandexHandbookML,prince2023understanding, ho2020denoisingdiffusionprobabilisticmodels, diffusionpolicy2023}, emphasis on \href{https://education.yandex.ru/handbook/ml/article/vvedenie-v-generativnoe-modelirovanie}{Yandex handbook} \citep{YandexHandbookML}, \href{https://udlbook.github.io/udlbook/}{UDL} book \citep{prince2023understanding}, and \href{https://www.youtube.com/watch?v=qJeaCHQ1k2w}{Deepai} video tutorials on GMs
\begin{itemize}
    \item \textbf{Theory:}
    Generative modeling goals: synthesis, compression, uncertainty, and \emph{planning via sampling}.
    VAEs: latent space as controllable representation (and typical failure modes).
     Generative modeling goals: synthesis, compression, uncertainty, and planning via sampling. VAEs: latent space as controllable representation (and typical failure modes). Diffusion models: denoising, conditioning, guidance; why they dominate image generation. Flow matching. Bridge to robotics: generative models for actions/trajectories (policy as a generative model). Foundation generative models overview: ChatGPT, Qwen, DALL·E, Midjourney, etc.
    \item \textbf{Skills:}
    Implement a VAE for reconstruction + latent interpolation.
    Generate diverse action candidates and re-rank them by a simple cost function.
    The use of known generative models in their inference mode, feel the joy of contemplation. 
\end{itemize}

\item \textbf{Vision-Language Models}
\citep{radford2021learningtransferablevisualmodels, Alayrac2022, blip2_2023, llava_2023, qwen_vl_2023, internvl_2024, zhang2024visionlanguagemodelsvisiontasks, hu2021loralowrankadaptationlarge, prince2023understanding}, epmhasis on basics of  \href{https://education.yandex.ru/handbook/ml/article/obuchenie-predstavlenij}{representation learning} \citep{YandexHandbookML} and \href{https://arxiv.org/abs/2103.00020}{CLIP} paper \citep{radford2021learningtransferablevisualmodels}
\begin{itemize}
    \item \textbf{Theory:}
    Why VLMs changed CV: open-vocabulary recognition, retrieval, grounding, and prompting as an API.
    Two dominant families:
    (1) \textbf{contrastive} alignment (CLIP-like) for retrieval/zero-shot,
    (2) \textbf{generative} multimodal LLMs (LLaVA/BLIP-2/Qwen-VL/InternVL-style) for reasoning and instruction following.
    Data and objectives: contrastive, captioning, instruction tuning, grounding.
    Practical adaptation: PEFT/LoRA, prompt engineering, and evaluation traps (leakage, overclaiming).
    \item \textbf{Skills:}
    Zero-shot + few-shot classification and retrieval with CLIP-like models.
    Fine-tune a VLM head with LoRA on a small domain dataset; write a compact evaluation report.
\end{itemize}

% ------------------------------------------------------------
\item \textbf{Segmentation and Object Detection}
\citep{prince2023understanding, ronneberger2015unet, 
ren2015faster, redmon2016yolo, carion2020detr, 
lin2017focal, rezatofighi2019giou,
he2017mask, kirillov2023segment, sam2_2024, groundingdino2023}, emphasis on \href{https://udlbook.github.io/udlbook/}{UDL} book \citep{prince2023understanding}, and original papers on \href{https://arxiv.org/abs/1505.04597}{U-net} \citep{ronneberger2015unet}, \href{https://arxiv.org/abs/1506.02640}{YOLO} \citep{redmon2016yolo}, \href{https://arxiv.org/abs/2304.02643}{SAM} \citep{kirillov2023segment}
\begin{itemize}

    \item \textbf{Theory:}
    Segmentation taxonomy (semantic: U-Net, DeepLab; instance: Mask R-CNN; panoptic). 
    Object detection taxonomy (two-stage: Faster R-CNN; one-stage: YOLO; transformer-based: DETR). 
    Architectural components: anchor-based vs anchor-free, Feature Pyramid Networks (FPN), NMS vs end-to-end set prediction (Hungarian matching). 
    Losses and metrics: IoU/GIoU/DIoU, focal loss, mAP@0.5 and mAP@[0.5:0.95], IoU/Dice for segmentation. 
    Promptable and open-vocabulary perception: SAM → SAM 2, language-grounded detection (GroundingDINO-style). 
    Model selection logic for robotics and real-time systems.

    \item \textbf{Skills:}
    Train and evaluate a YOLO detector on a custom dataset. 
    Fine-tune a segmentation model (U-Net or DeepLab) and compute IoU/Dice. 
    Compare YOLO, Faster R-CNN, and DETR in terms of speed and mAP. Run SAM (or SAM2) for automatic and prompt-based segmentation and evaluate mask quality. Build a pipeline combining detection and SAM refinement. Perform structured error analysis (small objects, occlusion, domain shift, imbalance).
\end{itemize}

% ------------------------------------------------------------
\item \textbf{Maps of Depth, Landmarks Detection}
\citep{YandexHandbookML, prince2023understanding, ming2021deep, kumar2020luvli}, emphasis on \href{https://huggingface.co/learn/computer-vision-course/en/unit8/3d_measurements_stereo_vision}{stereo vision} intuition \citep{HuggingFaceCVcourse}, \href{https://arxiv.org/abs/2406.09414}{Depth Anything} \citep{yang2024depthv2} and landmark detection \href{https://arxiv.org/abs/2004.02980}{LUVLi} \citep{kumar2020luvli} original papers, and \href{https://udlbook.github.io/udlbook/}{pose estimation} basics \citep{prince2023understanding}
\begin{itemize}
    \item \textbf{Theory:}
    Depth: stereo vs monocular; where each breaks; evaluation pitfalls (scale ambiguity).
    Keypoints/pose: heatmaps vs regression, uncertainty and occlusion handling.
    Practical robotics angle: 2D keypoints $\rightarrow$ 3D pose; calibration; tracking stability.
    \item \textbf{Skills:}
    Build a depth map (stereo or pretrained monocular) and validate on a small benchmark.
    Run pose/landmarks with a pretrained model (MediaPipe/dlib or modern equivalents), then stabilize temporally (simple filtering).
\end{itemize}

% ------------------------------------------------------------
\item \textcolor{gray}{\textbf{Midterm Exam}}

% ------------------------------------------------------------
\item \textbf{Face Recognition}
\citep{zeng2021survey, liu2023patch, turk1991eigenfaces, koch2015siamese, schroff2015facenet, deng2019arcface, wang2021deep, prince2023understanding}
\begin{itemize}
    \item \textbf{Theory:}
    Metric learning pipeline: embeddings, margins, sampling, and evaluation (verification vs identification).
    Robustness: pose, occlusion, aging, domain shift.
    Security/ethics: bias, privacy, deepfakes, and responsible deployment (what you can measure, what you cannot claim).
    \item \textbf{Skills:}
    Fine-tune a face embedding model (or use pretrained) and evaluate verification ROC / FAR-FRR trade-offs.
    Real-time detection+tracking+recognition pipeline (with clear limitations documented).
\end{itemize}


% ------------------------------------------------------------
\item \textbf{Plus 1D: Video and 3D Image Processing}
\citep{maturana2015voxnet, Wang_2017, mao2021voxeltransformer3dobject, prince2023understanding, singh20203d}
\begin{itemize}
    \item \textbf{Theory:}
    Video: temporal modeling (3D CNNs, video transformers), tracking, and action recognition.
    ``Foundation'' trend: tracking/segmentation across frames (prompting and memory, conceptually).
    3D: reconstruction basics (SfM/MVS) + modern view: neural rendering/3D representations as compressed scene models.
    \item \textbf{Skills:}
    Fine-tune a video model on a small action dataset or do self-supervised features + linear probe.
    Implement multi-object tracking (SORT/DeepSORT style) and evaluate on a short sequence.
    Build a small 3D reconstruction demo (Open3D) and document failure cases.
\end{itemize}

% ------------------------------------------------------------
\item \textbf{Cloud of Points Processing}
\citep{YandexHandbookML, qi2017pointnetdeeplearningpoint, lai2022stratified}
\begin{itemize}
    \item \textbf{Theory:}
    Representations: points vs voxels vs implicit fields; why point transformers work.
    Core models: PointNet family + graph/DGCNN intuition + transformer-style point processing.
    Robotics angle: LiDAR perception, registration, segmentation of drivable space/objects.
    \item \textbf{Skills:}
    Run PointNet/PointNet++ baseline; then compare with a transformer-like point model (if compute allows).
    Segment a small LiDAR scene or do object classification; report metrics + qualitative visualizations.
\end{itemize}

% ------------------------------------------------------------


% ------------------------------------------------------------
\item \textbf{Vision-Language-Action Models: Part \#1}
\citep{rt1_2022, rt2_2023, openx_2023, openvla_2024, octo_2024, diffusionpolicy2023}
\begin{itemize}
    \item \textbf{Theory:}
    What makes a VLA/VLAM different from a VLM: actions as a modality, closed-loop control, and safety constraints.
    Core design choices:
    (a) \textbf{action representation} (continuous vs discretized/action tokens),
    (b) \textbf{policy learning} (behavior cloning, offline RL, diffusion policies),
    (c) \textbf{conditioning} (language goals, goal images, histories/memory).
    Canonical lines:
    RT-1 (robot transformer scaling),
    RT-2 (web knowledge transfer to actions),
    Open X-Embodiment / RT-X (multi-robot data standardization),
    OpenVLA and Octo as open(-ish) generalist policies.
    Compute-realism constraints: why we often do ``small VLA'' in sim + evaluate carefully.
    \item \textbf{Skills (expanded):}
    Build a ``VLM-to-action prerequisite'' pipeline:
    CLIP retrieval + grounded segmentation (CLIPSeg/SAM-style) $\rightarrow$ object-centric state.
    Implement imitation learning baseline in simulation (or a toy manipulator) conditioned on language.
\end{itemize}

% ------------------------------------------------------------
\item \textbf{Vision-Language-Action Models: Part \#2}
\citep{saycan_2022, palme_2023, gemini_robotics_2025}
\begin{itemize}
    \item \textbf{Theory:}
    System architectures for embodied agents:
    \textbf{end-to-end} VLA vs \textbf{modular} (perception $\rightarrow$ planner $\rightarrow$ skills).
    Planning + grounding: SayCan-style ``LLM proposes, skills constrain''; embodied reasoning models (PaLM-E-style).
    Safety and deployment: constraint layers, uncertainty, guardrails, and the need for measurable safety claims.
    Real-world constraints: latency, on-device inference, calibration, logging, data flywheel (collect $\rightarrow$ fine-tune $\rightarrow$ evaluate).
    The ``next trend'' to watch: robotics foundation models and agentic tool use in physical tasks.
    \item \textbf{Skills (expanded):}
    Mini-VLAM capstone:
    language instruction $\rightarrow$ grounded perception $\rightarrow$ symbolic plan (skills) $\rightarrow$ execution in sim/ROS2.
    Evaluate with a task suite: success rate + failure taxonomy + short demo video.
\end{itemize}

% ------------------------------------------------------------
\item \textbf{A lecture of the student's choice}
\begin{itemize}
    \item \textbf{Approaching Artificial General Intelligence (AGI)} \citep{deepseekai2025deepseekr1incentivizingreasoningcapability}
    \item \textbf{Dynamic Visual Reasoning by Learning Differentiable Physics} \citep{}
    \item \textbf{Info-structure and Deploy} \citep{}
    \item \textbf{Mixing and Tuning of the Models} \citep{}
    \item \textbf{Transformer from scratch step-by-step} \citep{}
\end{itemize}

\end{enumerate}

% ============================
\section{Practical sessions}
% ============================
\begin{enumerate}
    \item \textbf{Feature Extraction and Machine Learning} \citep{gitCV-2026, LightningColab}
    \item \textbf{Convolutional Neural Networks} \citep{gitCV-2026}
    \item \textbf{Visual Transformers}  \citep{gitCV-2026}
    \item \textcolor{black}{Generative Models \citep{gitCV-2026}}%{\textbf{Assignment \# 01}}
    \item \textbf{Vision-Language Models and Zero (and/or Few) - Shot Classification} \cite{gitCV-2026}
    \item \textbf{Segmentation and/or Object Detection} \citep{gitCV-2026}
    \item \textbf{Maps of Depth and/or Pose Estimation} \citep{gitCV-2026}
    \item \textcolor{gray}{\textbf{Midterm Exam}}
    \item \textbf{Face recognition and Object Tracking} \citep{gitCV-2026}
    \item \textbf{3D Data Reconstruction and/or Video Processing} \citep{gitCV-2026}
    \item \textbf{Cloud of Points Processing} \cite{gitCV-2026, YandexHandbookML}
    \item \textbf{Vision-Language-Action Models: Part \#1} \citep{gitCV-2026}
    \item \textbf{Vision-Language-Action Models: Part \#2} \citep{gitCV-2026}
    \item \textbf{Deploy of a CV model} \citep{gitCV-2026}
    \item \textcolor{gray}{\textbf{Defend of the students' projects}}
\end{enumerate}

% ============================
\section{Group Project Track (assignment \#1), week 02}
% ============================
The course project must involve \textbf{visual data (pixels)} or \textbf{multimodal data} (vision+text+sensor).
Projects are encouraged to align with \textbf{Robotics and Computer Vision} and emphasize strong evaluation and reproducibility.

\subsection*{Suggested project topics (choose one)}
\textbf{VLM-focused (perception + prompting):}
\begin{enumerate}
\item \textbf{Open-vocabulary detector for lab objects:} Grounded detection + promptable segmentation; evaluate on novel categories and lighting changes.
\item \textbf{CLIP-based anomaly detection:} learn ``normal'' with self-supervision, use text prompts to explain anomalies; evaluate on a small industrial/robotics dataset.
\item \textbf{Domain adaptation with LoRA:} adapt a VLM to a niche domain (medical/robotics parts) with 200--1000 labeled pairs; compare PEFT vs full fine-tuning.
\item \textbf{Video understanding with VLMs:} retrieve key moments from robot videos via text queries; measure retrieval quality and latency.
\item \textbf{Referring segmentation benchmark:} build a mini-dataset of referring expressions in your lab; evaluate grounded detection + segmentation baselines.

\end{enumerate}

\textbf{VLA/VLAM-focused (capstone, recommended):}
\begin{enumerate}
\setcounter{enumi}{5}
\item \textbf{Mini-VLAM in simulation:} language instruction $\rightarrow$ grounded perception $\rightarrow$ skill execution; report success rate and failure taxonomy.
\item \textbf{Diffusion policy for manipulation:} implement a diffusion-based action generator in sim; compare vs behavior cloning.
\item \textbf{SayCan-style planner:} LLM proposes a sequence of skills, skills constrain feasibility; evaluate on a kitchen/tabletop task set.
\item \textbf{Open X-Embodiment style data format:} unify 2--3 small robot dataset
\end{enumerate}

% ============================
\section{Competition Track (assignment \#2, week 09)}
% ============================
\textbf{A Computer Vision Competition with Kaggle Private Dataset}
\begin{itemize}
    \item \textbf{Challenge (as an example):} ``Open-vocabulary perception for robotics scenes'' (recognize + localize novel objects from text prompts).
    \item \textbf{Task:} open-vocabulary detection/segmentation + a simple downstream metric (e.g., pick-list correctness).
    \item \textbf{Dataset:} private lab/robot camera images + hidden test split; prompts provided for evaluation.
    \item \textbf{Skills:} baselines (CNN/ViT), strong models (grounded detection + SAM-style segmentation), experiment tracking, report writing.
\end{itemize}

\section{Midterm}
60 min test. 
\section{Final Exam}
60 min test.

\bibliography{references}
\end{document}